\subsection{de Sitter Vacua in String Theory --- arXiv:hep-th/0301240}

\textbf{Authors:} Shamit Kachru, Renata Kallosh, Andrei Linde, Sandip P. Trivedi

\paragraph{主な主張}
Type IIBフラックスコンパクト化で、非摂動的超ポテンシャルによる超対称AdS真空を反D3ブレーンでupliftし、モジュライを固定した準安定de Sitter真空を構成する機構を提案した。 \cite{Kachru:2003aw}

\paragraph{新規性}
フラックスによる固定とDブレーン・インスタントンまたはゲージーノ凝縮を組み合わせ、残るK\"ahlerモジュライの固定と正の真空エネルギーを同時に扱った。

\paragraph{位置付け}
KKLTシナリオの原論文であり、弦理論におけるde Sitter真空構成とその制御可能性の議論の出発点である。

\paragraph{コメント（読書メモ）}
読書時の分類は「modular sym.なし」（この表現の適用範囲は未確認）。元リストの文献[1]に当たるGKP構成\cite{Giddings:2001yu}へ反D3ブレーンを導入する文脈で読み、KKLTの原論文として整理する。
